\documentclass[a4paper,12pt,titlepage]{jsarticle}

\usepackage{otf}

\usepackage{graphicx}
\usepackage{verbatim}
\usepackage{amsmath}
\usepackage{listings}
\usepackage{array}


\pagestyle{plain}

\title{プログラミングB レポート}
\author{課題番号：2\\提出年月日：\today\\提出者：09B25036 HU XIANDONG\\教員名：武政 淳二}
\date{}

\begin{document}

\maketitle

%%%%%%%%%%%%%%%%%%%%%%%%%%%%%%%%%%%%%%%%%%%%
\section{課題内容}
%%%%%%%%%%%%%%%%%%%%%%%%%%%%%%%%%%%%%%%%%%%%

本レポートでは，「課題２」に取り組んだ。
具体的には，中置記法で与えられた算術式に対して，
再帰下降構文解析を用いて四則演算・括弧・指数演算を解析し，
値を計算する電卓プログラムを作成した。
発展課題として，小数の扱い，指数演算，一元負号の処理，
および文法エラー検出（発展1, 2a, 2b, 3, 4）を実装した。

%%%%%%%%%%%%%%%%%%%%%%%%%%%%%%%%%%%%%%%%%%%%
\section{アルゴリズムの説明}
%%%%%%%%%%%%%%%%%%%%%%%%%%%%%%%%%%%%%%%%%%%%

本課題では，中置記法の式を次の文法に基づいて解析する
再帰下降構文解析法を用いた。

\begin{itemize}
    \item $expression ::= term\ \lbrace\ (+ \mid -)\ term\ \rbrace$
    \item $term ::= factor\ \lbrace\ (* \mid /)\ factor\ \rbrace$
    \item $factor ::= constant \mid (expression) \mid -factor \mid factor^{\text{digit}}$
    \item $constant ::= digits\ [.\ digits]$
\end{itemize}

式を「項」，項を「因子」に分解し，
各規則を関数に対応させ，
左から順に入力文字を読み進めながら再帰的に評価する。

アルゴリズムの要点は次の通りである。

\begin{itemize}
    \item 全体の計算は $expression$ を入口とする。
    \item $term$ が乗除を，$factor$ が括弧・指数・一元負号を処理する。
    \item 文字列の走査には 1 文字の前読み（lookahead）を用い，
    関数 \verb|next()| により現在文字 \verb|current_char| を更新する。
    \item 文法に合致しない文字を検出した場合は，入力位置とともに
    エラーを出力して終了する。
\end{itemize}

%%%%%%%%%%%%%%%%%%%%%%%%%%%%%%%%%%%%%%%%%%%%
\section{プログラムの説明}
%%%%%%%%%%%%%%%%%%%%%%%%%%%%%%%%%%%%%%%%%%%%

\subsection{入力形式}
本プログラムは，標準入力から 1 行の算術式を読み込む。
空白の混入は想定しない。

\subsection{出力形式}
式が正しい場合は計算結果を実数で出力し，
除算におけるゼロ除算では
\verb|Zero division occurs.| と表示し終了する。
文法エラーの場合は，
\verb|Syntax error occurs at the X-th character.| と出力する。

\subsection{データ構造}
入力式は文字配列 \verb|input[]| に格納し，
現在位置を整数 \verb|char_index|，
現在の文字を \verb|current_char| として保持する。

\subsection{大域変数}
\begin{verbatim}
char input[MAX_EXPRESSION_LENGTH];
int char_index;
char current_char;
\end{verbatim}

\subsection{関数の設計}

関数間の呼び出し関係は図\ref{fig:callgraph}の通りである。

\begin{figure}[h]
\centering
\begin{verbatim}
                    +-------------+
                    | expression  |
                    +------+------+
                           |
                           v
                    +-------------+
                    |    term     |
                    +------+------+
                           |
                           v
                    +-------------+
                    |   factor    |
                    +------+------+
                           |
                           v
                    +-------------+
                    |  constant   |
                    +-------------+
\end{verbatim}
\caption{関数呼び出し関係図}
\label{fig:callgraph} 
\end{figure}

\subsection{関数の説明}
本節では，各関数がどのように文法規則に対応し，
入力をどのように消費するかを説明する。

\paragraph{next()}
入力文字列 \verb|input| の走査位置 \verb|char_index| を 1 進め，
現在の文字 \verb|current_char| を更新する関数である。
本プログラムは 1 文字先読み方式を採用しているため，
構文解析の各段階で \verb|next()| を呼び出すことで字句を順次消費する。

\paragraph{syntax\_error()}
文法規則に反した入力が検出された際に呼び出される関数である。
現在位置 \verb|char_index| を用いて，
「何文字目でエラーが発生したか」を出力し，
プログラムを強制終了する。
構文解析中の例外処理として機能する。

\paragraph{constant()}
連続する数字列を読み取り，整数部および小数部から実数値を生成する関数である。
まず整数部を構成する数字を読み取り，続いて小数点が存在する場合は
小数部の各桁を 0.1 倍しながら加算する。
この関数は文法規則
\[
constant ::= digits\ [\,.\, digits\,]
\]
を実装している。

\paragraph{factor()}
最も優先順位の高い構文要素を扱う関数である。
この関数は以下の 4 種類の要素に対応する：

\begin{itemize}
    \item 数値（\verb|constant()| の呼び出し）
    \item 括弧で囲まれた式（\verb|(| に続く \verb|expression()| の再帰呼び出し）
    \item 一元負号（\verb|-factor()| による符号反転）
    \item 指数演算（\verb|factor^{digit}|）
\end{itemize}

特に括弧を含む場合，\verb|(| を読み飛ばした後に \verb|expression()| を呼び出し，
対応する \verb|)| が存在することを確認することでネスト構造を処理する。

\paragraph{term()}
\verb|factor| を単位として，乗算および除算を処理する関数である。
文法規則は以下の通りである：

\[
term ::= factor\ \lbrace\ (* \mid /)\ factor\ \rbrace
\]

先頭で \verb|factor()| を読み取り，
その後に続く \verb|*| または \verb|/| を検出した場合は，
右側の \verb|factor| を評価し，演算を逐次適用する。
除算において右辺が 0 の場合には，
課題仕様に従い \verb|Zero division occurs.| を出力して終了する。

\paragraph{expression()}
式全体を評価する最上位の関数であり，加算および減算を処理する。
文法規則は以下の通りである：

\[
expression ::= term\ \lbrace\ (+ \mid -)\ term\ \rbrace
\]

まず先頭の \verb|term| を読み取り，
その後に続く \verb|+| または \verb|-| を検出するたびに
後続の \verb|term| を再帰的に取得して値を更新する。
本関数の戻り値が式全体の評価結果となる。

\newpage
\appendix
\section{プログラム}
\begin{verbatim}
   1: #include <stdio.h>
   2: #include <stdlib.h>
   3: #include <ctype.h>
   4: #include <string.h>
   5:
   6: #define MAX_EXPRESSION_LENGTH 100
   7:
   8: char input[MAX_EXPRESSION_LENGTH];
   9: int char_index = 0;
  10: char current_char;
  11:
  12: void next();
  13: void syntax_error();
  14: double constant();
  15: double factor();
  16: double term();
  17: double expression();
  18:
  19: int main(void) {
  20:     fgets(input, sizeof(input), stdin);
  21:
  22:     int len = strlen(input);
  23:     if (input[len - 1] == '\n')
  24:     {
  25:         input[len - 1] = '\0';
  26:     }
  27:
  28:     char_index = 0;
  29:     next();
  30:
  31:     double result = expression();
  32:
  33:     if (current_char != '\0') {
  34:         syntax_error();
  35:     }
  36:
  37:     printf("%.10g\n", result);
  38:
  39:     return 0;
  40: }
  41:
  42: void next() {
  43:     current_char = input[char_index++];
  44: }
  45:
  46: void syntax_error() {
  47:     printf("Syntax error occurs at the %d-th character.\n", char_index);
  48:     exit(0);
  49: }
  50:
  51: double constant() {
  52:
  53:     if (!isdigit(current_char))
  54:     {
  55:         syntax_error();
  56:     }
  57:
  58:     double value = 0;
  59:
  60:     while (isdigit(current_char))
  61:     {
  62:         value = 10 * value + (current_char - '0');
  63:         next();
  64:     }
  65:
  66:     if (current_char == '.')
  67:     {
  68:         next();
  69:
  70:         if (!isdigit(current_char))
  71:         {
  72:             syntax_error();
  73:         }
  74:
  75:         double decimal_scale = 0.1;
  76:
  77:         while (isdigit(current_char))
  78:         {
  79:             value += (current_char - '0') * decimal_scale;
  80:             decimal_scale *= 0.1;
  81:             next();
  82:         }
  83:     }
  84:
  85:     return value;
  86: }
  87:
  88: double factor() {
  89:     double value = 0;
  90:
  91:     if (current_char == '-')
  92:     {
  93:         next();
  94:         return -factor();
  95:     }
  96:
  97:     if (current_char == '(')
  98:     {
  99:         next();    // skip'('
 100:         value = expression();
 101:
 102:         if (current_char != ')')
 103:         {
 104:             syntax_error();
 105:         }
 106:         next();     // skip')'
 107:     }
 108:
 109:     else if (isdigit(current_char))
 110:     {
 111:         value = constant();
 112:     }
 113:
 114:     else {
 115:         syntax_error();
 116:     }
 117:
 118:     if (current_char == '^')
 119:     {
 120:         next();
 121:
 122:         if (!isdigit(current_char))
 123:         {
 124:             syntax_error();
 125:         }
 126:
 127:         int exponent = current_char - '0';  // one-digit exponent required
 128:         next();
 129:
 130:         double result = 1;
 131:         for (int i = 0; i < exponent; i++)
 132:         {
 133:             result *= value;
 134:         }
 135:
 136:         return result;
 137:     }
 138:
 139:     return value;
 140: }
 141:
 142: double term() {
 143:     double value = factor();
 144:
 145:     while (current_char == '*' || current_char == '/')
 146:     {
 147:         char op = current_char;
 148:         next();
 149:
 150:         double right_operand = factor();
 151:
 152:         if (op == '*')
 153:         {
 154:             value *= right_operand;
 155:         }
 156:         else {
 157:
 158:             if (right_operand == 0)
 159:             {
 160:                 printf("Zero division occurs.\n");
 161:                 exit(0);
 162:             }
 163:
 164:             value /= right_operand;
 165:         }
 166:     }
 167:
 168:     return value;
 169: }
 170:
 171: double expression() {
 172:     double value = term();
 173:
 174:     while (current_char == '+' || current_char == '-')
 175:     {
 176:         char op = current_char;
 177:         next();     // skip + or -
 178:
 179:         double right = term();
 180:
 181:         if (op == '+')
 182:         {
 183:             value += right;
 184:         }
 185:         else {
 186:             value -= right;
 187:         }
 188:     }
 189:
 190:     return value;
 191: }
\end{verbatim}

\end{document}